\documentclass[12pt,a4paper]{article}
\usepackage[utf8]{inputenc}
\usepackage[T2A]{fontenc}
\usepackage[russian]{babel}
\usepackage{amsmath,amssymb,amsthm}
\usepackage{booktabs,array}
\usepackage{microtype}
\usepackage{geometry}
\geometry{margin=2.3cm}
\setlength{\emergencystretch}{2em}

\newtheorem{proposition}{Предложение}
\newtheorem{nogocriterion}{Критерий провала}

\title{Проверка гипотезы меню состояний\\
Пост-S2T-II.A exploratory gate}
\author{}
\date{4 августа 2026 г.}

\begin{document}
\maketitle

\section{Гипотеза}

Рассмотрим намеренно радикальную интерпретацию:
\begin{quote}
Компактный объект
\[
K=\mathbb{RP}^3\times S^1
\]
является не дополнительным физическим пространством, а пространством меток глобальной склейки квантовых состояний. Частица должна пониматься как устойчивый цикл или сектор этого меню.
\end{quote}

Эта гипотеза рассматривается как новая версия III.0 и не используется для возвращения закрытых утверждений II.A.

\section{Минимальное меню}

Фундаментальная группа носителя равна
\[
\pi_1(K)=\mathbb Z_2\times\mathbb Z.
\]
Поэтому пространство плоских \(U(1)\)-характеров имеет вид
\[
\mathcal M_{flat}
=\operatorname{Hom}(\mathbb Z_2\times\mathbb Z,U(1))
\cong\{+1,-1\}\times U(1).
\]
Таким образом, минимальное пространство состояний состоит из двух несвязных
компонент, каждая из которых изоморфна \(U(1)\). Координата на \(U(1)\)
задаёт непрерывную фазу, а номер компоненты --- дискретную торсионную ветвь.

Spin-структуры образуют torsor над
\[
H^1(K,\mathbb Z_2)\cong\mathbb Z_2\oplus\mathbb Z_2,
\]
поэтому их четыре. Если считать spin-структуру частью меню, расширенное пространство состоит из восьми окружностей. Если spin-структура является фиксированным фоном, физически доступными остаются только две окружности.

\begin{proposition}[Топологические метки устойчивых циклов]
Свободные homotopy-классы петель в \(K\) маркируются парой
\[
(\varepsilon,n)\in\mathbb Z_2\times\mathbb Z.
\]
Поэтому базовая топология естественно даёт один целый winding charge и один дискретный parity label.
\end{proposition}

\section{Первый положительный результат}

Интерпретация \(K\) как меню математически непротиворечива. Она даже естественнее прежней буквальной compactification-интерпретации в одном отношении: объёмы, spin-структуры и голономии читаются как характеристики пространства допустимых глобальных состояний, а не обязаны непосредственно задавать микроскопические пространственные размеры.

Гипотеза тем самым проходит первый gate:
\begin{center}
\begin{tabular}{>{\raggedright\arraybackslash}p{0.35\textwidth}>{\raggedright\arraybackslash}p{0.50\textwidth}}
\toprule
Проверка & Результат \\
\midrule
Дискретная фазовая ветвь & да, \(\mathbb Z_2\) \\
Циклическая фазовая координата & да, \(U(1)\) \\
Целый winding label & да, \(n\in\mathbb Z\) \\
Spin-sector labels & да, четыре global choices \\
\bottomrule
\end{tabular}
\end{center}

\section{Неабелев no-go}

Однако
\[
\mathbb Z_2\times\mathbb Z
\]
является абелевой группой. Любое конечномерное унитарное представление её генераторов состоит из коммутирующих нормальных матриц. Такие матрицы одновременно диагонализуются, а представление распадается на одномерные характеры.

\begin{nogocriterion}[Провал полного particle menu]
Базовые петли \(K\) не могут сами по себе породить неприводимый \(SU(2)\)-дублет или \(SU(3)\)-триплет. Неабелевы weak- и color-мультиплеты требуют дополнительного fiber, groupoid или operator algebra над меню состояний.
\end{nogocriterion}

Это можно увидеть на простом примере. Два генератора \(\pi_1(K)\) обязаны коммутировать. Матрицы Паули \(\sigma_x\) и \(\sigma_z\), напротив, имеют ненулевой коммутатор и потому не могут представлять эти две петли одновременно.

\section{Что меню не объясняет}

Топология классифицирует сектора, но не задаёт:
\begin{itemize}
    \item энергию или массу каждого сектора;
    \item вероятность перехода между секторами;
    \item три поколения --- естественные числа меню равны \(2\), \(4\) и \(8\), но не \(3\);
    \item неабелевы gauge charges;
    \item локальную динамику и пространство-время.
\end{itemize}

Следовательно, гипотеза проходит как кинематический скелет, но закрывается отрицательно как самостоятельная теория частиц.

\section{Следующий gate}

Минимальное допустимое расширение имеет вид
\[
\begin{gathered}
\text{две окружности состояний}
+\text{единый неабелев fiber}
\\
+\text{фиксированный оператор переходов}
\end{gathered}
\]
Следующая проверка должна ответить на один конкретный вопрос: существует ли один минимальный неабелев fiber, который без независимых секторных настроек порождает одновременно \(SU(2)\)-дублет и \(SU(3)\)-триплет. До такого вывода меню состояний нельзя использовать для вычисления масс или физических констант.

\section{Вердикт}

Первая сумасбродная гипотеза не оказалась бессмысленной. Она даёт ясное пространство глобальных меток: две фазовые окружности, четыре spin-структуры и заряды \((\varepsilon,n)\). Но она также немедленно обнаруживает собственную границу: базовая топология полностью абелева. Поэтому идея перспективна как язык кинематической классификации, но не как готовая модель микромира.

\end{document}